% Generated draft: requires mathematical and source review.
\documentclass[11pt]{article}
\usepackage[margin=25mm]{geometry}
\usepackage{fontspec}
\setmainfont{Latin Modern Roman}
\usepackage{amsmath,amssymb,mathtools,unicode-math}
\setmathfont{Latin Modern Math}
\usepackage{enumitem,needspace,xurl,hyperref}
\setlength{\parindent}{0pt}
\setlength{\parskip}{5pt}
\setlength{\emergencystretch}{3em}
\newcommand{\N}{\mathbb{N}}
\newcommand{\Z}{\mathbb{Z}}
\newcommand{\Q}{\mathbb{Q}}
\newcommand{\R}{\mathbb{R}}
\newcommand{\C}{\mathbb{C}}
\newcommand{\F}{\mathbb{F}}
\newcommand{\A}{\mathbb{A}}
\newcommand{\PP}{\mathbb{P}}
\newcommand{\E}{\mathbb{E}}
\newcommand{\eps}{\varepsilon}
\newcommand{\dd}{\,\mathrm{d}}
\begin{document}
\section*{UnsolvedMath ID: 30001507}
\textbf{Source problem code:} catalogue entry
\subsection{Definitions and mathematical statement}
A group \(\Gamma\) is \emph{coherent} if every finitely generated subgroup of \(\Gamma\) is finitely presented. Thus \(\Gamma\) is \emph{noncoherent} if there exists a subgroup \(H\leq\Gamma\) that is finitely generated but not finitely presented. A lattice in a Lie group \(G\) is a discrete subgroup \(\Gamma<G\) for which \(G/\Gamma\) has finite invariant volume. Two Lie groups are locally isomorphic when their Lie algebras are isomorphic.

\textbf{Original source formulation, as supplied.} The record gives the following conjecture: if \(G\) is a semisimple Lie group not locally isomorphic to \(\operatorname{SL}_2(\R)\) or \(\operatorname{SL}_2(\C)\), then every lattice in \(G\) is noncoherent. Equivalently, the asserted conclusion for each such lattice \(\Gamma<G\) is
\[
\exists H\leq\Gamma\quad H\text{ is finitely generated and is not finitely presented}.
\]
The original formulation does not explicitly specify connectedness, whether the Lie groups are real, or restrictions concerning compact groups, compact factors, or the center. These qualifications require source review; none is added here.

\textbf{Conflicting normalized formulation.} The fields \texttt{statement} and \texttt{clean\_statement} instead assume that \(G\) is a connected semisimple real Lie group, with the same two local-isomorphism exclusions, and ask whether every infinite-index subgroup of a lattice \(\Gamma<G\) has infinite covolume. If covolume is interpreted in the ambient group \(G\), this asks
\[
\forall H\leq\Gamma,\qquad [\Gamma:H]=\infty
\ \Longrightarrow\ \operatorname{vol}(G/H)=\infty.
\]
Here \([\Gamma:H]\) is the number of cosets of \(H\) in \(\Gamma\), and volume means invariant quotient volume induced by a fixed Haar measure on \(G\). The choice of ambient group for covolume is not explicit in the supplied question. This formulation concerns volume and index, whereas the original concerns finite presentation of finitely generated subgroups. The record supplies no equivalence between them, so they are retained as distinct formulations pending review. Neither the intended target nor the missing hypotheses should be inferred from the title alone.
\subsection{Short English statement}
The supplied original conjecture says that every lattice in the specified semisimple Lie groups contains a finitely generated subgroup that is not finitely presented. The normalized question instead concerns infinite index and covolume. The intended statement and its full hypotheses require human source review. The supplied open label is a curation claim, not an independently verified status.
\subsection{Sources}
\begin{itemize}
\item Local registry: supplied UnsolvedMath record, permanent ID \(30001507\), source problem code \texttt{catalogue entry}, in the Oberwolfach Reports Open Problems collection. The record is titled ``Noncoherence of Higher-Rank Semisimple Lattices.'' No individual proposer is supplied; the recorded proposed year is 2010.
\item Supplied report citation: \emph{Geometric Group Theory, Hyperbolic Dynamics and Symplectic Geometry} (2011), Oberwolfach Reports, pp.\ 1761--1825, DOI \texttt{10.4171/owr/2010/30}. \url{https://doi.org/10.4171/owr/2010/30}. The background describes these reports as documenting mathematical workshops and their open-problem discussions.
\item Supplied literature reference labeled ``Kapovich noncoherence theorem'': \url{https://doi.org/10.2140/gt.2013.17.39}. No further bibliographic details are supplied.
\item Supplied literature reference labeled ``original OWR report'': \url{https://doi.org/10.4171/owr/2010/30}.
\item Literature assessment dated exactly \texttt{2026-08-21}: this is supplied dated curation, not independently verified. It labels the problem open and describes the record as a fragment referring to a broad noncoherence conjecture for lattices. It reports that a 2013 primary account proves important arithmetic cases while leaving general nonarithmetic rank-one lattices open. The background explicitly calls this dated literature triage and requests independent verification of time-sensitive status and cited claims.
\item The supplied record labels its normalized statement \texttt{corrected\_verified} and says its mathematical content was checked against the official report, with notation, context, and question form normalized without changing the source-level claim. This is a supplied verification assertion; the conflict with the supplied original formulation remains unresolved here. No external source has been visited for this draft.
\end{itemize}
\end{document}
